\section{Physical origin of finite-interface transport bias}
\subsection{Hydraulic response of a diffuse viscosity profile}
Consider two equal-density fluids layered across a no-slip slit of width $H$,
driven parallel to their interface by uniform force $F$. In steady low-speed
flow, $d(\mu du/dy)/dy=-F$. Define the inverse-viscosity moments
\begin{equation}
 I_n=\int_0^H\frac{y^n}{\mu(y)}\,dy,\qquad n=0,1,2.
\end{equation}
Integrating the momentum balance and imposing $u(0)=u(H)=0$ gives
\begin{equation}
 \mu u'=F(I_1/I_0-y),\qquad
 Q=\int_0^H u\,dy=F\left(I_2-\frac{I_1^2}{I_0}\right).
 \label{eq:conductance}
\end{equation}
This elementary continuum relation evaluates the response of a specified
constitutive profile without solving the lattice dynamics. A sharp step,
the viscosity reconstructed from measured phase fractions, and a separately
calibrated tanh profile therefore provide distinct references. At unit
equal reference density, dynamic and kinematic viscosities coincide to
leading order; small density departures remain recorded in the data.

Let $s=y-H/2$, $\varphi=(1-\tanh(s/\xi))/2$, and
$r=\mu_1/\mu_0$. For a well-separated interface, the integrated excess shear
compliance relative to the sharp step is
\begin{equation}
 \mathcal C_A=\frac{\xi}{2}\left(\frac1{\mu_1}-\frac1{\mu_0}\right)\ln r,
 \qquad \mathcal C_H=0.
\end{equation}
The subscripts denote arithmetic and harmonic interpolation. Harmonic
interpolation cancels this zeroth-moment excess, but the stress varies across
the interface in pressure-driven flow. Its first and second weighted moments
still change the conductance. Writing $\gamma=(r-1)/(r+1)$, expansion of
Eq.~\ref{eq:conductance} about the centered sharp interface gives
\begin{align}
 \frac{Q_A-Q_{sharp}}{Q_{sharp}}&=
 -\frac{3\gamma^3\ln r}{4-3\gamma^2}\frac{\xi}{H}
 +O((\xi/H)^2),\\
 \frac{Q_H-Q_{sharp}}{Q_{sharp}}&=
 \frac{2\pi^2\gamma^2}{4-3\gamma^2}\left(\frac{\xi}{H}\right)^2
 +O((\xi/H)^4).
 \label{eq:thininterface}
\end{align}
Exponential wall-tail corrections vanish in the separated-interface limit.
The supplement derives these expressions. They describe the stated
one-dimensional constitutive problem, not universal error laws for curved
interfaces or heterogeneous rock. Harmonic interpolation's improved leading
order therefore does not imply superior accuracy at every confinement.

Three independent unforced planar calculations calibrate $W_{0.9}$ at
$\beta=0.35,0.7,0.9$. The widths are 9.376, 4.765 and 3.929 nodes;
$\xi=W_{0.9}/[2\operatorname{atanh}(0.9)]$ specifies a tanh approximation
without fitting any velocity. Eighteen forced calculations combine these
parameters with $H=16,32,64$, both closures and $(\tau_0,\tau_1)=(0.6,1.4)$.
Force scales as $H^{-3}$ and surface tension as $H^{-1}$, with fixed
relaxation times. Fixed $\beta$ makes the physical interface shrink during
this diffusive-scale sequence, so it is not a fixed-interface spatial
convergence test. All late relative velocity changes are below $10^{-11}$.

At $H=16$, $\beta=0.35$, arithmetic interpolation decreases the sampled
volume flux by 9.19\%, while harmonic interpolation increases it by 21.59\%.
Independent planar-tanh predictions give $-9.08$\% and $+21.30$\%.
The corresponding profile discrepancies against the \emph{sharp} solution
are 17.02\% and 24.66\%, but discrepancies against the continuum solution
using linear interpolation of the measured diffuse profile are 0.160\% and
0.100\%. Shape-preserving cubic interpolation gives 0.121\% and 0.273\%.
Most of the apparent failure is consequently the response of the chosen
finite constitutive transition, rather than failure of the solver to recover
its continuum momentum balance.

Fig.~\ref{fig:transport} also shows the competing effect of sharpening an
interface. At $H=16$, increasing $\beta$ from 0.35 to 0.9 reduces the
arithmetic sharp-flux bias, but increases its measured-diffuse profile error
from 0.160\% to 0.753\% with the specified linear reconstruction; the
corresponding shape-preserving cubic values are 0.121\% and 0.443\%.
At $H=64$, $\beta=0.9$, the harmonic discrepancy is 0.0117\% for linear
interpolation and 0.0482\% for the cubic reference. Quadrature refinement
changes a reference profile by at most $9.12\times10^{-8}$ relative to the
measured velocity, whereas interpolation changes can be much larger.
These small residuals include reference-reconstruction sensitivity and
finite-density departures; they are not isolated lattice truncation errors.
Interface confinement and lattice representation
must therefore be assessed separately. The remaining error is not assigned
exclusively to one discrete operator, and these flat interfaces do not test
curved-interface spurious currents or capillary trapping.

\begin{figure}[htbp]\centering
\includegraphics[width=.95\linewidth]{interfacial-transport.pdf}
\caption{Controlled transport at viscosity ratio 9:1. Left: measured flux
shifts against the midpoint-sampled sharp profile; lines show continuous
tanh-profile predictions from the independent width calibration, with no
flow-data fitting. Right: velocity-profile error against the continuum
reference reconstructed from measured phase fractions. Bars span the
linear/PCHIP reconstructions and quadrature refinement, not statistical
confidence intervals. Diffuse constitutive bias and a reconstruction-dependent
solver discrepancy are distinct observables.\AIPlotDisclosure}
\label{fig:transport}\end{figure}

\subsection{Wall direction and capillary-stress magnitude}
The capillary second moment is proportional to
$|\vct g|(\vct n\vct n-I)$, so setting $\vct n$ alone does not set the stress
magnitude. On a flat D2Q9 lower wall, a locally linear gradient has the
clipped-stencil transformation
\begin{equation}
 \vct g_{clip}=\operatorname{diag}(5/6,1/2)\vct g,\qquad
 \frac{|\vct g_{clip}|}{|\vct g|}
 =\sqrt{\frac{25}{36}\sin^2\theta+\frac14\cos^2\theta}.
 \label{eq:wallmagnitude}
\end{equation}
The recovered wall rotation preserves the magnitude of the gradient it
receives. Even after an imposed direction correction, the local perturbation
amplitude thus retains this linear-field attenuation: 0.764 at 60 or 120
degrees and 0.833 at 90 degrees. Moment reconstruction removes that specific
linear-stencil bias. It does not prove discrete mechanical equilibrium at a
nonlinear contact line or prescribe an accurate curved-wall normal. The
measured static-angle improvement and increased parasitic flow reported
below are consistent with needing both checks; Eq.~\ref{eq:wallmagnitude}
alone is not a quantitative prediction of the observed capillary pressure
or spurious velocity.

\section{Numerical controls for transport and capillarity}
All runs use double precision, GCC 13.3, CMake 3.28.3 and Open MPI 4.1.6 on
Linux. Dedicated workstation timings are described separately below.
The tested physical range is low-speed flow at unit, equal reference density,
$\tau_k\in\{0.6,0.8,1,1.4\}$ and prescribed $\sigma=0.005$ for droplet cases.
The supplement supplies derivations, measurement protocols and run identities.

\subsection{Force, viscosity and restart}
In the original MRT operator, $S$ was fixed at the first-phase relaxation
time. Pure second-phase pressure-channel fields were therefore identical
when its relaxation time changed. The original velocity-shift force has
momentum increment $s_j\tau F$: D3Q19 conserved rates make it vanish, while
D2Q9's old conserved rates produce a relaxation-dependent force.
The same fixed relaxation also changes capillary stress: at unchanged
prescribed $\sigma=0.005$, swapping $(\tau_0,\tau_1)=(0.6,1.4)$ to
$(1.4,0.6)$ gives inferred original surface tensions 0.00315 and 0.00764.
Local relaxation restores the intended surface-tension scale.

Sixteen uniform-force cases spanning both models, lattices, collision
operators and two viscosities agree with $u(100)=10^{-4}$ within
$3.2\times10^{-16}$. Twenty-four force-driven channel cases use halfway walls
and compare with $u(y)=F_xy(H-y)/(2\mu)$. Force scales as $H^{-3}$ while
relaxation time is fixed, preserving diffusive physical scaling. SRT errors
fall by approximately four on each doubling from $H=8$ to 16 to 32.
MRT errors are about $1.4\times10^{-8}$ or smaller, consistent with the
remaining transient at the selected run duration (Fig.~\ref{fig:channel}).

\begin{figure}[htbp]\centering
\includegraphics[width=.94\linewidth]{channel-verification.pdf}
\caption{Analytic halfway-wall channel verification. SRT refinement and MRT
profiles are measured independently from globally reconstructed VTI output.\AIPlotDisclosure}
\label{fig:channel}\end{figure}

Continuous and restarted uniform-force and three-dimensional droplet runs
have bitwise-identical macroscopic fields in six comparisons, including
one-to-two-rank restart. Invalid parameter cases exit with failure. Algebraic
checks of the emitted production coefficients have maximum absolute error
below $9\times10^{-16}$.

For two equal-density layers at viscosity ratio 9:1, the arithmetic
interpolation has 5.37\% relative profile error at $H=64$ against the
sharp-interface solution. Its error against a continuum solution using the
measured diffuse viscosity profile is 0.22\%. Harmonic interpolation reduces
the sharp-interface discrepancy to 0.81\%, while the diffuse-reference error
is 0.23\% (Fig.~\ref{fig:layered}). Thus the major change is a constitutive
transition bias, not an improved collision formula.

\begin{figure}[htbp]\centering
\includegraphics[width=.94\linewidth]{layered-flow.pdf}
\caption{Layered Poiseuille flow: sharp-interface and measured diffuse-profile
references diagnose viscosity-interpolation bias. The small residual also
depends on reference reconstruction.\AIPlotDisclosure}
\label{fig:layered}\end{figure}

\subsection{Interfaces and wetting}
Planar interfaces conserve component mass and have spurious speed below
$5\times10^{-16}$. The measured distance between $\chi=-0.9$ and $+0.9$ is
4.765 nodes for D2Q9 and 5.289 for D3Q19 at $\beta=0.7$; D2Q9 widths at
$\beta=0.35$ and 0.9 are 9.376 and 3.929 nodes respectively. Interface width
is measured rather than equated to the input parameter.

Pressure is averaged in bulk regions outside the diffuse interface, and
radius is obtained by orthogonal fitting of subgrid contour intersections.
At 12000 steps, D2Q9 Laplace errors for initial radii 12, 20 and 28 are
1.50\%, 0.50\% and 0.23\%. The D3Q19 errors for radii 8, 12 and 16 are
4.96\%, 1.92\% and 0.99\%; the largest case is continued to 8000 steps.
The smallest spheres contain few interface widths and quantify finite-width
bias, rather than establish a sharp-interface asymptotic regime
(Fig.~\ref{fig:laplace}). Unequal-relaxation cases also retain the prescribed
surface-tension scale; their detailed errors are tabulated in the supplement.

\begin{figure}[htbp]\centering
\includegraphics[width=.94\linewidth]{laplace-verification.pdf}
\caption{Laplace pressure compared with $\sigma/R$ and $2\sigma/R$.
Insets quantify finite-interface error; late-time drift is checked separately.\AIPlotDisclosure}
\label{fig:laplace}\end{figure}

Sessile-drop measurements use the physical halfway wall $y=0.5$ and measure
the angle through phase 1. A principal orthogonal circle fit excludes points
within $0.75W_{0.9}$ of the wall. Sensitivities to contour level and exclusion
height, left--right asymmetry and late-time changes are retained.
For initial radius 36 and $\beta=0.7$, the equilibrated 60-degree measurements
are 64.52 degrees for the clipped stencil and 60.91 degrees for
renormalization. The corresponding 120-degree measurements are 117.18 and
120.92 degrees. Neutral-angle bias persists, and renormalization increases
its spurious speed from approximately $2.8\times10^{-5}$ to $1.05\times10^{-4}$.
This is a tradeoff rather than uniform superiority (Fig.~\ref{fig:wetting}).

\begin{figure}[htbp]\centering
\includegraphics[width=.94\linewidth]{wetting-sensitivity.pdf}
\caption{Flat-wall contact-angle and spurious-velocity ablation. Error envelopes
show deterministic measurement sensitivity rather than statistical confidence.
Size and interface-width sequences use fixed lattice surface tension and are
not labelled a matched-physics spatial convergence study.\AIPlotDisclosure}
\label{fig:wetting}\end{figure}

\subsection{Direct closed-slit capillary response}
Closed periodic slits contain a central phase-0 segment and two menisci, with
halfway-wall spacing $H$, length $4H$ and no external force. Pressure is
measured directly in separated bulk buffers, rather than inferred from
velocity. The signed sharp reference is
$p_0-p_1=2\sigma\cos\theta/H$, with the angle through phase 1.
The H16 pilot tests 60, 90 and 120 degrees, three gradients and both capillary
operators. The follow-up uses $H=16,32,48$ for legacy perturbation,
linear-reconstructed perturbation and linear-reconstructed CSF.
At fixed $\beta=0.7$ and $\tau_0=\tau_1=1$, $\sigma=0.005/s$ and the initial
duration is $8000s^2$, where $s=H/16$. The three-width 60-degree controls
are continued to twice that duration. This is a joint grid/sharp-interface sequence, not
spatial convergence at fixed physical interface width.

% Generated from capillary-response.json; hold-change bound rounded upward.
\newcommand{\CapillaryRenPressureErrorFine}{1.655}
\newcommand{\CapillaryCSFPressureErrorFine}{1.714}
\newcommand{\CapillaryNonneutralReductionFine}{3.539}
\newcommand{\CapillaryNeutralReductionFine}{9.109}
\newcommand{\CapillaryHoldPressureChange}{0.004}

For the 60-degree controls, the legacy signed pressure discrepancy remains
approximately $-6.3$ to $-7.5$\% over these three widths. Linear
reconstruction reduces its discrepancy to \CapillaryRenPressureErrorFine\%
at $H=48$; CSF gives \CapillaryCSFPressureErrorFine\%.
CSF lowers the H48 nonneutral residual maximum speed by a factor
\CapillaryNonneutralReductionFine\ relative to reconstructed perturbation,
while the neutral-wall factor is \CapillaryNeutralReductionFine\
(Fig.~\ref{fig:closedcapillary}). Lower velocity therefore does not imply
lower pressure bias. The longest-hold checks change normalized pressure by
less than \CapillaryHoldPressureChange\% of the sharp reference, with the
late color RMS changes retained separately. The data do not establish a
nonzero asymptotic legacy error from three widths or an entry threshold.
The largest longer-hold color RMS change is $4.82\times10^{-9}$, and the
maximum relative component/total mass drift across those earlier slit histories is
$4.03\times10^{-12}$.
Neutral residual $\mu U_{max}/\sigma$ remains near $3.0\times10^{-3}$ for
linear-reconstructed perturbation and $3.3\times10^{-4}$ for CSF over the
tested widths. Decreasing lattice velocity under diffusive scaling therefore
does not demonstrate vanishing dimensionless parasitic flow. This is a
specific limitation for low-rate capillary pore-flow studies, not a universal
error floor for other geometries or wetting schemes.

The negative controls are informative. Local tanh reconstruction increases
the H16 neutral maximum speed from $9.05\times10^{-5}$ to
$9.79\times10^{-5}$ relative to linear reconstruction with perturbation.
Setting $\sigma=0$ leaves a wall-to-center color-profile difference of about
0.0109 while fluid speed falls below $2\times10^{-16}$. Thus color-profile
distortion occurs independently of capillary motion. The evidence motivates
inspection of segregation and boundary reflection together; it does not
prove that one discrete term exclusively causes the residual flow.
Bulk CSF Laplace controls retain errors of 1.39\% and 0.303\% for D2Q9
radii 12 and 24 after longer relaxation, and 4.26\% for D3Q19 radius 8.
Those tests do not validate curved three-dimensional wetting.

\begin{figure}[htbp]\centering
\includegraphics[width=.98\linewidth]{closed-capillary.pdf}
\caption{Closed-slit pressure and residual velocity. Left: signed 60-degree
pressure discrepancy after the longer holds; bars retain bulk-sampling and
duration sensitivity. Right: $\mu U_{max}/\sigma$, with filled circles for
60 degrees and open triangles for neutral wetting. Methods share the input
surface tension and measured bulk width. The joint refinement sequence
reduces pressure error for reconstruction, while CSF reduces residual flow.
It does not remove the wall-color distortion.\AIPlotDisclosure}
\label{fig:closedcapillary}\end{figure}

\subsection{Wall phase transport, equilibrium and physical location}
A neutral zero-sigma control isolates the phase operator from capillary
feedback. At rest, its axial scalar recurrence is
\begin{equation}\label{eq:phaserecurrence}
 \varphi'_j-\varphi_j=A\Delta_x\varphi_j+
 \beta B(J_{j-1}-J_{j+1}),\qquad
 J_j=n_j\varphi_j(1-\varphi_j).
\end{equation}
For cosine recoloring, $A=a+2b$ and $B=a+\sqrt2b$, whereas local
wall bounce-back gives $A_w=a+b$ and $B_w=a+b/\sqrt2$;
$a=1/9$ and $b=1/36$. The wall/bulk stationary ratios differ.
A cosine-specific tangent multiplier
$k=[(a+b)(a+\sqrt2b)/(a+2b)-b/\sqrt2]/a=0.9511844635$
restores that ratio. It does not restore the entire transient operator.
Published projection recoloring already has equal stationary ratios,
so this multiplier is not needed for that model.
Independent scalar recurrences, cutoff/rotation controls and the
unfavorable nonneutral results are detailed in the supplement.

% Generated from wall-consistency.json; bounds checked and rounded upward.
\newcommand{\WallZeroColorBound}{$9\times10^{-16}$}
\newcommand{\WallFixedColorBound}{$2\times10^{-13}$}
\newcommand{\WallFixedSpeedBound}{$7\times10^{-16}$}
\newcommand{\WallRawSpeed}{$5.780\times10^{-8}$}
\newcommand{\WallPairSpeedBound}{$1\times10^{-14}$}
\newcommand{\WallStaggeredEnergyPercent}{99.9999999}

The cosine correction preserves generated neutral equilibria within
\WallFixedColorBound\ in color and \WallFixedSpeedBound\ in velocity.
From a step initialization, however, raw adjacent fields can retain
period-two momentum (Fig.~\ref{fig:wallbalance}). A spatially alternating
equilibrium momentum perturbation is preserved by collision and reverses
under streaming, giving the established eigenvalue $-1$
\cite{lallemand2000,guo2011}. Even-time reporting can conceal it.
Adjacent-step stopping checks reject such false stationarity; no filtering
or damping is applied to numerical fields. Exploratory nonneutral use
of the cosine multiplier worsens the measured pressure bias.

\begin{figure}[htbp]\centering
\includegraphics[width=.98\linewidth]{wall-balance.pdf}
\caption{Neutral discrete fixed point and temporal controls. A cosine-specific
wall correction removes static profile distortion but can retain raw
alternating velocity. Adjacent-step averages diagnose this mode;
they are not substituted into the solver. Complete histories and
nonneutral negative controls are supplied in the supplement.\AIPlotDisclosure}
\label{fig:wallbalance}\end{figure}

The full transient instead requires a conservative face transfer. For
adjacent wall nodes $L,R$ and outgoing reflected diagonals $p,m$, define
\begin{equation}\label{eq:wallphaseflux}
 J_{LR}=f_{0p}(L)-f_{0m}(R)
 -\frac{\varphi_L+\varphi_R}{2}[f_{Tp}(L)-f_{Tm}(R)].
\end{equation}
Opposite incoming face populations receive $\pm J_{LR}$; the second phase
is complementary to its original total reflection. Each component is
conserved globally, total-fluid reflection is unchanged and pure phases
have zero transfer. At unit-density rest this reduces to the established
forward specular scalar transport \cite{ginzburg2015}. It restores the
complete neutral bulk recurrence, with six full-solver trajectories
agreeing within \FluxTransientBound\ and unfiltered neutral velocity below
\FluxNeutralSpeedBound. This condition controls numerical phase reflection,
not a general contact-line constitutive relation.

Physical wall location remains distinct. Rotating the normal at the first
fluid center imposes the angle half a cell from the halfway wall. A circular
sharp-interface estimate then replaces $H$ by $H-1$, giving pressure bias
$1/(H-1)$; this is a testable geometric estimate rather than a fitted width.
The halfway control extrapolates the inward normal to the physical wall.
The published scalar ghost formula \cite{liu2015} instead sets
$\chi_g=\chi_1-\cot\theta|1.5g_{t1}-0.5g_{t2}|$, followed by a full
virtual-color stencil without a second center-angle rotation.
This formula alone is not the complete published algorithm.

Phase transfer with ordinary center rotation worsens pressure. With ghost
color, the 60-degree errors after four initial durations become
\FluxGhostPressureSixteen\%, \FluxGhostPressureThirtyTwo\% and
\FluxGhostPressureFortyEight\% at H16/H32/H48 (Fig.~\ref{fig:wallflux}).
Nonneutral currents, contour-fitting sensitivity and adjacent-step phase
oscillations remain. The matched-core comparison below therefore retains
signed pressure errors and raw velocities. Neutral transient consistency
and static pressure improvement establish neither nonneutral dynamic
accuracy nor curved digital-rock wetting.

\begin{figure}[htbp]\centering
\includegraphics[width=.98\linewidth]{wall-flux.pdf}
\caption{Wall transport and location are separate controls. (a) Independent
phase-transient discrepancy. (b) Raw neutral velocity. (c) Signed nonneutral
pressure discrepancy; the grey curve is the first-center geometric estimate.
(d) Nonneutral residual capillary numbers. The full neutral transient is
preserved without implying negligible nonneutral circulation.\AIPlotDisclosure}
\label{fig:wallflux}\end{figure}

\subsection{Published core controls at matched phase transport}
The ghost-color comparison initially retains the recovered recoloring and
MRT definitions. To separate these from wall treatment, two independent
controls reproduce the explicitly specified core operators of
Liu et al. \cite{liu2015}. Projection recoloring uses
\begin{equation}
 R_i=\beta_p\rho\varphi(1-\varphi)w_i(\vct e_i\cdot\vct n),
\end{equation}
instead of dividing this projection by $|\vct e_i|$.
The published D2Q9 raw-moment spectrum, in the existing row order, is
$(1,1.63,1.14,1,1.92,1,1.92,1/\tau,1/\tau)$.
The matching Guo source and physical half-force velocity preserve the
momentum increment even when its relaxation rate is one.
Only CSF is admitted with this spectrum: its energy rates do not justify
reusing the legacy perturbation-tension mapping. The fluid-only
normal derivative at walls remains an explicit LBFlow choice; these
controls are not a full replication of an incompletely specified scheme.

For an axial rest profile, projection segregation has $B_p=A=a+2b$,
and local bounce-back has $B_{p,w}=A_w=a+b$. It already preserves the
bulk stationary ratio without the cosine-derived multiplier.
To compare transport independently of pressure, set
\begin{equation}
 \beta_p=\beta_c\frac{a+\sqrt{2}b}{a+2b}
             =\CoreMatchedBeta\quad(\beta_c=0.7).
\end{equation}
This matches the complete axial bulk recurrence, not just a fitted width.
Independent 128-cell planar profiles agree within
$\CoreMatchedProfileError$; both have $W_{0.9}=4.765235$.
Projection with the unadjusted $\beta_p=0.7$ instead has width 4.450814.
The cosine weights $w_i/|\vct e_i|$ have a fourth-order isotropy defect
$T_{xxxx}-3T_{xxyy}=0.0650874$, whereas standard projection weights
give zero to roundoff. Equal axial recurrence therefore does not make
the multidimensional discrete operators identical. The supplied oblique
width diagnostic also retains its different sampling interval and
interpolation bias; its full difference is not attributed to this tensor.

With ghost color and phase transfer, the matched projection/legacy-MRT
pressure discrepancies are $+0.517$, $-0.336$ and $-0.702$\% at
$H=16,32,48$ after four initial durations. The published-MRT values are
$+0.551$, $-0.341$ and $-0.716$\% (Fig.~\ref{fig:core}).
The sign reversal and larger final discrepancy preclude using the close
H16 result as uniform accuracy evidence. Raw nonneutral speeds decrease,
but adjacent color still differs from its two-step state. At H16 the
one-step change is approximately $6.66\times10^{-7}$, compared with
$3.47\times10^{-12}$ over two steps. Projection and a different kinetic
spectrum do not eliminate the conserved-momentum lattice mode.

Its force coupling can be checked exactly. On an even periodic axis, or
the reflecting wall-normal axis, let $P_t$ be the fluid-node mean of
physical momentum times the alternating spatial sign, and $G_t$ the same
projection of force density. Outgoing momentum is $P_t+G_t/2$; streaming
or reflection changes its sign. Hence
\begin{equation}\label{eq:parityforce}
 P_{t+1}=-P_t+\frac{G_{t+1}-G_t}{2}.
\end{equation}
Physical fields, native outgoing populations and the actual force audit
satisfy this relation within $5.1\times10^{-19}$ in the original held
nonneutral control. Explicitly setting $\sigma=0$ makes circulation
decay while alternating momentum remains; this diagnostic changes the
physics and does not filter production velocities. Odd-periodic-domain
and exact-translation controls are distinguished in the supplement.

The MRT change also affects wall hydraulics. In a pure low-speed slit,
define $\Lambda=(\tau-1/2)(1/s_q-1/2)$ and
$C=4\Lambda/3-1/4$. The midpoint profile is
\begin{equation}
 u(y)=\frac{F}{2\nu}[y(H-y)+C],\qquad y=j+\tfrac12.
\end{equation}
Legacy rates give $\Lambda=3/16$ and $C=0$;
the published $s_q=1.92$ gives $\Lambda=(\tau-1/2)/48$.
Eighteen pure-channel controls match this independent reference within
relative profile error $\CoreMaxHydraulicError$.
The largest unshifted flux discrepancy is approximately $-2.30$\% at
$H=8$, $\tau=0.6$. A lower static current cannot be interpreted as
unchanged dynamic wall response. Bulk droplet controls at matched
transport likewise retain their measured pressure and residual velocity.

\begin{figure}[htbp]\centering
\includegraphics[width=.98\linewidth]{published-core.pdf}
\caption{Independently matched core controls. (a) Axial planar profiles;
matched projection overlaps cosine, while equal beta changes the width.
(b) Signed 60-degree pressure discrepancies after longer holds.
(c) Corresponding unfiltered residual capillary numbers.
(d) Pure-channel flux offsets of the published MRT spectrum: symbols are
measurements and lines are the independent midpoint prediction.\AIPlotDisclosure}
\label{fig:core}\end{figure}

\subsection{Guarded pressure-driven fronts}
A pressure boundary introduces component mass, so closed-mass conservation
alone cannot validate it. The new experiment admits two opposed axial PB
planes only with flat ghost-color walls and phase transfer. At every
incoming/outgoing state, each prescribed phase must remain pure within
$10^{-9}$ on its PB plane and two interior strips. The actual first-strip
gradient must be below the existing $10^{-9}$ cutoff, keeping the
uncorrected Zou--He reconstruction force-free. The solver fails when the
interface reaches this collar. Phase-only wall transfers on faces touching
PB nodes are disabled, including the wrap between opposite reservoirs.
Velocity reservoirs, films and irregular walls are outside this extension.

For interior VD/GV nodes, independently integrate every outgoing kinetic
population entering from or leaving for PB. Solid links return by reflection
and are excluded. The exact accounting is
\begin{equation}
 M_k(t)-M_k(t_0)=\sum_{n=t_0}^{t-1}
   [\mathcal J_{k,\mathrm{PB}\to\mathrm{int}}(n)
    -\mathcal J_{k,\mathrm{int}\to\mathrm{PB}}(n)]+\epsilon_k.
\end{equation}
The boundary-node sum of $\rho u$ is not substituted for these link fluxes.
Restart establishes a new accounting baseline while retaining the same
outgoing numerical state. Two-phase fields and populations are bitwise
identical across the tested rank/restart changes. Across the pressure
campaign, the largest budget residual relative to initial interior mass
is $\CoreMaxPressureBudget$.

Pure-channel pressure controls use both flow directions, both phases and
both spectra, at $H=8,16,32$ and $\tau=1,1.55$. The independently measured
central pressure gradient and specified midpoint offset reproduce the
central profiles within $1.92\times10^{-5}$. The central gradient differs
slightly from the nominal endpoint pressure divided by length, so measured
pure-channel conductance is retained separately.

A recorded crop of the held closed meniscus produces a single front in
a $4H$ slit. Remote corridors explicitly combine donor phase populations
into the prescribed pure phase, preserving each total population and its
momentum. No damping or surface-tension fit is applied. Equal viscosities,
60-degree walls and independently matched beta are used at H16/H32;
$\sigma$ and duration scale as in the closed controls.
The volume-equivalent front coordinate is
$X_V=1/2+H^{-1}\sum_{\mathrm{int}}\varphi$; centerline and first-fluid-row
crossings are retained as separate diagnostics.

Reservoir differences $0.98,1,1.02$ times
$p_Y=2\sigma\cos\theta/H$ give both motion directions. Increasing the
duration fourfold to 32000/128000 steps changes the three-point linear
zero-drift pressure ratio by at most $\CoreMaxDurationRatioChange$.
Additional pressure pairs, prespecified at the separately measured closed
pressure plus/minus $0.00025p_Y$, still give opposite drift signs at
both durations. The linear zero-drift estimate differs from the closed
pressure by at most \CoreMaxZeroClosedPercent\% of $p_Y$
(Fig.~\ref{fig:fronts}; supplementary table). This is a finite-time flat-slit
consistency check, not a geometric pore-throat trapping threshold.

The measured slope of front speed against pressure remains approximately
11.3--17.8\% lower than pure-channel conductance. A zero-sigma diagnostic at unchanged reservoir
pressure recovers the pure-channel mean velocity within relative
$6.1\times10^{-6}$ and front speed within $5.9\times10^{-5}$,
supporting total-fluid decoupling rather than interpreting phase segregation
alone as hydraulic resistance. Contact-line dissipation and its
regularization are established physical issues \cite{huh1971,akai2018};
these comparisons do not isolate a physical slip length, a constitutive
contact-line law or a general numerical-error mechanism. Static pressure
agreement alone is insufficient evidence for quantitative displacement
rates in digital-rock pores.

\begin{figure}[htbp]\centering
\includegraphics[width=.98\linewidth]{pressure-fronts.pdf}
\caption{Pressure-driven single-front controls after four durations.
(a,b) Both drift directions at H16/H32.
(c) Estimated zero-drift pressure and independently measured closed pressure.
(d) Front mobility relative to pure-channel conductance. Ghost color and
phase transfer are shared; projection beta is independently matched.
The narrower prespecified sign brackets and all raw mass budgets remain
in the supplied data. No trapped pore or hysteresis loop is simulated.\AIPlotDisclosure}
\label{fig:fronts}\end{figure}

\subsection{Fixed interface width and localized front resistance}
The preceding fixed-beta sequence also sharpens the interface in physical
units. To separate that change, projection recoloring and legacy MRT are
retained at $\tau_0=\tau_1=1$, $\theta=60^\circ$ and
$W_{0.9}/H=4.765235021846117/16$. Beta is selected from the independent
descending axial zero-flux condition
\begin{equation}
 \varphi_{j+1}-\varphi_j=-\beta[
 \varphi_j(1-\varphi_j)+\varphi_{j+1}(1-\varphi_{j+1})].
\end{equation}
Production planar controls verify the resulting widths at H16/H32/H48
to relative $\FrontWidthRelativeError$. Beta is approximately
$0.63165825,0.31145624,0.20699672$; no pressure or velocity enters that
calibration. Diffusive units retain the leading axial scalar diffusion and
fluid viscosity, both $1/6$, as well as the measured relative width.
They do not independently fix an effective contact-line slip length.

Closed donors supply pressure centers before any moving-front run.
Offsets $\pm0.02p_Y$ are used on all three grids, with additional
$\pm0.01p_Y,\pm0.04p_Y$ pairs and lengths $4H,8H$ at H16.
Doubling the front duration changes the measured mobility by at most
\FrontDurationMobilityPercent\%. With $M_V=d\dot X_V/d\Delta p$ and
independently measured pure-channel mobility $G_0$, define
\begin{equation}
 R_e=M_V^{-1}-G_0^{-1}.
 \label{eq:front-resistance}
\end{equation}
The H16 mobility deficit decreases from 11.26\% to 5.92\% when the
channel doubles in length, while $R_e$ changes by only
\FrontLengthResistancePercent\%. The response is almost independent of
the fourfold pressure-halfwidth variation in this tested low-speed interval.
Thus a length-independent interfacial contribution, added to the
length-dependent hydraulic resistance, explains the measured length effect
(Fig.~\ref{fig:front-resistance}). This does not separately locate all
dissipation at the microscopic contact line.

The low-speed reciprocal balance provides an independent force diagnostic.
Let $\bm u_a,Q_a$ be the velocity and flux of a pure-channel auxiliary
driven by $\Delta p_a$. In incompressible creeping flow with the same
viscosity and stationary walls, the reciprocal work relation is
$\Delta p_a Q=\Delta p Q_a+\int\bm F_c\cdot\bm u_a\,dV$.
Its lattice quadrature defines
\begin{equation}
 p_F=-\frac{\sum_{\mathrm{fluid}}\bm F_c\cdot\bm u_a}{Q_a},
 \qquad \overline U\simeq G_0(\Delta p-p_F).
 \label{eq:force-work}
\end{equation}
The auxiliary is measured in separate positive/negative pure-channel runs;
the production capillary force is audited from two adjacent checkpoints.
Their reconstructed densities and color match physical output within
$6.7\times10^{-16}$. Two-state averaging is applied only to diagnostics.
The incremental resistance $\delta p_F/\delta\overline U$ agrees with
$\delta\Delta p/\delta\overline U-G_0^{-1}$ within
\FrontForceResistancePercent\%, with no force or slip parameter fitted.
Absolute speed residuals up to $\FrontMaxForceResidual$ remain, and can be
larger than the drift near rest. Equation~\eqref{eq:force-work} is consequently
not asserted as an exact discrete identity or a zero-drift estimator.

At fixed measured relative width the $4H$ mobility deficits are
11.26\%, 18.57\% and \FrontFineDeficit\% on the three grids.
Signed closed-pressure discrepancies are $+0.517$\%, $+1.389$\% and
$+1.692$\%. The response has not converged on these resolutions.
Force-work closure identifies how the numerical momentum balance produces
the deficit; it does not verify its magnitude against independent physical
contact-line data. The remaining discretization, diffuse regularization and
wall coupling require a separate, matched moving-contact-line reference
before quantitative rock displacement or pore-throat trapping is reported.

\begin{figure}[htbp]\centering
\includegraphics[width=.98\linewidth]{front-resistance.pdf}
\caption{Fixed-width and localized-resistance controls with unchanged
projection/legacy MRT. (a) Independently calibrated planar profiles in common
length units. (b) Grid-dependent mobility deficit at fixed measured $W_{0.9}/H$.
(c) Nearly length-independent excess resistance across the tested pressure
interval. The vertical scale resolves its small measured variation.
(d) Resistance from independently weighted capillary force versus pressure
and instantaneous mean flow, with the identity line for comparison.
All results retain raw flow and exact component budgets.\AIPlotDisclosure}
\label{fig:front-resistance}\end{figure}

